\chapter{Hilfe und Beratung}

\section{Erste Hilfe: GAF}

Wir kennen nicht immer die Lösung, wissen dafür aber meistens, wer sie
kennt. Wir haben gute Kontakte zu vielen Institutionen und Personen an
dieser Uni. Wenn du uns einfach mal besuchen willst, bist du herzlich
willkommen. (Kontakt siehe Kapitel \ref{gafKontakt}, S. \pageref{gafKontakt})


\section{Webforen und Kommilitonen}

In den Foren kannst du dich mit deinen Kommilitonen
(und teils auch mit Lehrpersonal) austauschen. Wenn du Fragen direkt zu den
Übungen oder der Vorlesung hast, kannst du dich auch einfach an die
Übungsleiter der jeweiligen Vorlesung wenden. Keine Sorge, die beißen nur
selten.

Wenn du im IRC (ein Chatroom) unterwegs bist, findest du unter \#gaf und \#informatik.lmu auf
freenode auch immer andere Studenten aus deinem Fach.

\begin{urlList}
	\urlItem{http://die-informatiker.net}
% 	\urlItem{http://die-physiker.org}
%	\urlItem{http://die-mathematiker.net}
\end{urlList}

\skiptobottom
\begin{center}
	{\includegraphics[height=3.5cm]{img/wisdom_of_the_ancients}}
\end{center}

%\begin{figure}[h]
%\centering
%\includegraphics[height=3cm]{img/wisdom_of_the_ancients}
%\caption{}
%\label{fig:wisdom_of_the_ancients}
%\end{figure}


\newpage %% Das soll wieder raus

\section{Probleme mit Lehrveranstaltungen oder Lehrpersonal}

Die offizielle Ansprechperson hierbei ist der Studiendekan deiner
Fakultät. Er ist für die Qualität der Lehre verantwortlich. In jedem
Fall ist der sinnvollste Weg zu einer Lösung erst einmal das direkte
Gespräch mit dem Dozenten. Erst wenn ihr das Gefühl habt, ein Problem
lässt sich nicht anders lösen, bittet euren Studiendekan um
Hilfe. Oder fragt uns von der GAF.

\subsection*{Studiendekane Fakultät 16}
Prof. Werner Bley (Mathematik)\subjectList{\subjectM\subjectW}\\
Prof. Hans Jürgen Ohlbach (Informatik)\subjectList{\subjectI\subjectMI}\\
Prof. Thomas Augustin (Statistik)\subjectList{\subjectS}

\subsection*{Studiendekan Fakultät 17}
Prof. Dr. Jochen Weller \subjectList{\subjectP}

\section{Ansprechpartner nach Studiengängen}

Alle nachfolgenden Personen sind sehr umgängliche Menschen, mit denen
man bestens reden kann. Wie die meisten Professoren beißen sie nicht,
wenn man etwas zu beanstanden hat.

\subsection*{Mathematik (B.Sc., LA Gymnasium)\subjectList{\subjectM}}
PD Dr. Heribert Zenk (\mail{Heribert.Zenk@mathematik.uni-muenchen.de})\\
Theresienstraße 39, B333, Telefon: 089 / 2180 \emd{} 4460

PD Dr. Hartmut Weiß (\mail{hartmut.weiss@mathematik.uni-muenchen.de})\\
Theresienstraße 39, B317, Telefon: 089 / 2180 \emd{} 4680\\
Sprechstunde: Do, 15:00\emd{}16:00~Uhr

\subsection*{Wirtschaftsmathematik (B.Sc.)\subjectList{\subjectW}}
Prof. Dr. Gregor Svindland (\mail{studienberatung.wirtschaftsmathematik@math.lmu.de})\\
Theresienstraße 39, B231, Telefon: 809 / 2180 \emd{} 4628

\subsection*{Mathematik (LA Grund-, Haupt-, und Realschule)\subjectList{\subjectM}}
Dr. Erwin Schörner (\mail{schoerner@lmu.de})\\
Theresienstraße 39, B237, Telefon: 089 / 2180 \emd{} 4498

\subsection*{Mathematik (Fachdidaktik und Didaktik)\subjectList{\subjectM}}
Prof. Dr. Hedwig Gasteiger (\mail{gasteiger@math.lmu.de})\\ %TODO
Theresienstraße 39, B215, Telefon: 089 / 2180 \emd{} 4631

\subsection*{Informatik (B.Sc.)\subjectList{\subjectI}}
Dr. Reinhold Letz (\mail{reinhold.letz@lmu.de})\\
Oettingenstraße 67, E001, Telefon: 089 / 2180 \emd{} 9693\\
Sprechstunde: Di \& Mi 13:00\emd{}14:00~Uhr und nach Vereinbarung

\subsection*{Informatik (LA)\subjectList{\subjectI}}
Prof. Martin Hofmann, Ph.D. (\mail{lehramt@ifi.lmu.de})\\
Oettingenstraße 67, L107, Telefon: 089 / 2180 \emd{} 9341

\subsection*{Medieninformatik (B.Sc.)\subjectList{\subjectMI}}
Max Maurer / Simon Stusak (\mail{studentenbetreuer@medien.ifi.lmu.de})\\
Amalienstraße 17, 505, Telefon: 089 / 2180 \emd{} 4654

\subsection*{Physik (B.Sc.)\subjectList{\subjectP}}
Michael Rebhan (\mail{Michael.Rebhan@physik.uni-muenchen.de})\\
Schellingstraße 4, H417, Telefon: 089 / 2180 \emd{} 5033

\subsection*{Physik plus Meteorologie (B.Sc.)\subjectList{\subjectP}}
Dipl. Met. Heinz Lösslein (\mail{loesslein@lmu.de})\\
Theresienstraße 37, A208, Telefon: 089 / 2180 \emd{} 4217

\subsection*{Physik (LA)\subjectList{\subjectP}}
Prof. Dr. Raimund Girwidz (\mail{girwidz@physik.uni-muenchen.de})\\
Theresienstraße 37, A012, Telefon: 089 / 2180 \emd{} 2020


\section{Prüfungsamt}
Die Prüfungsämter sind für alle Prüfungsangelegenheiten zuständig,
also unter anderem für deine Noten, deine Praktika, deine Notenübersichten und
Abschlusszeugnisse. Sie sind bei der Fakultät zu finden, zu der
dein Studienfach gehört. Eine Zuordnung der Prüfungsämter zu den
einzelnen Studiengängen/-fächern findest du auf der Übersichtsseite
Studiengänge A--Z am unteren Ende der jeweiligen
Studiengangsinformationen.

\begin{urlList}
	\urlItem{http://www.lmu.de/pruefungsaemter}
\end{urlList}


\section{Studentenkanzlei}

Die Studentenkanzlei muss wegen gewissen formalen Belangen
gelegentlich besucht werden. Der Besuch dieses kafkaesken Molochs ist
oft mit großen Wartezeiten und Unbill verbunden. Es hilft, hartnäckig
zu bleiben und notfalls mehrfach zu kommen, bis du den richtigen
Sachbearbeiter triffst. Nicht umgehen lässt sich ein Besuch bei:

\begin{itemize}
\item Beantragen von Beurlaubungen (Krankheit, Ausland, Kinder\ldots)
\item Fragen zur Studienplatzvergabe/\,Immatrikulation (Anerkennung von Hochschulzugangsberechtigungen, nachträgliches Einschreiben, Verlust der Immatrikulationsbescheinigung)
\item Studienfachwechsel, zusätzliche Einschreibung für ein Doppelstudium
\item Bescheinigungen für die Krankenkasse und Rente, Quittungen für Studienbeiträge
\end{itemize}

Die Studentenkanzlei ist in den Räumen E011 und E114 im Hautgebäude zu finden und Montag bis Mittwoch, sowie am Freitag von 8:30 bis 11:30~Uhr geöffnet, sowie am Donnerstag von 13:30 bis 15:00~Uhr.

\begin{urlList}
	\urlItem{http://www.lmu.de/studentenkanzlei}
\end{urlList}


\section{Studieren mit Kind}

Auch für Eltern ist Studieren nicht unmöglich. Die Uni bietet diverse Beratungs"~ und Betreuungsmöglichkeiten.

\begin{urlList}
	\urlItem{http://studentenwerk-muenchen.de/studieren-mit-kind}
	\urlItem{http://www.lmu.de/studium/beratung/beratung_service/beratung_lmu/schwangere_kind}
\end{urlList}

%\section{Gleichstellungsbeauftragte}
%
%\begin{urlList}
%	\urlItem{www.gleichstellungsbeauftragte.lmu.de}
%\end{urlList}

\section{Die Frauenbeauftragten}
Weitere Anlaufstellen im Uni-Alltag vor allem bei Fragen und Problemen bezüglich Diskriminierungen und sexueller Belästigung im Wissenschaftsbetrieb sind die Frauenbeauftragten.
Das Aufgabengebiet der Frauenbeauftragten ist vielfältig und groß, darum hat zusätzlich zur Universitätsfrauenbeauftragten jede Fakultät eigene Frauenbeauftragte.

Alle Studierende können an den Weiterbildungsprogramm LMU-PLUS, welches durch das Büro der Frauenbeauftragten organisiert und aus Studienersatzmitteln finanziert wird, teilnehmen. Ausschließlich zur Förderung von Frauen ist das LMUMentoring und die Beratung zur finanziellen Förderung von Nachwuchswissenschaftlerinnen gedacht.

\begin{urlList}
	\urlItem{http://www.frauenbeauftragte.uni-muenchen.de/weiterbildung/plus/index.html}
	\urlItem{http://www.mathematik-informatik-statistik.uni-muenchen.de/fakultaet/beauftragte/index.html}
	\urlItem{http://www.physik.uni-muenchen.de/fakultaet/einrichtungen/frauenbauftragte/index.html}
\end{urlList}


\section{Studieren mit Behinderung}

Solltest du aufgrund einer Behinderung mehr Zeit, spezielle Hilfsmittel oder einen eigenen Raum für Klausuren benötigen, so kannst du beim Prüfungsamt einen Nachteilsausgleich beantragen.

\begin{urlList}
	\urlItem{http://studentenwerk-muenchen.de/studieren-mit-behinderung}
	\urlItem{http://www.lmu.de/barrierefrei}
\end{urlList}


\section{Student und Arbeitsmarkt}

Der Career Service der Universität bietet dir eine Stellen- und Praktikavermittlung, Kompetenztrainings, ein Mentoringprogramm, verschiedene Recruitung Events und einiges mehr. Einen Überblick verschaffst du dir am besten online oder du besuchst ihn in Ludwigstraße 27 / I. Stock am Montag, Dienstag, Donnerstag und Freitag zwischen 10 und 12~Uhr.

\begin{urlList}
	\urlItem{http://www.s-a.lmu.de}
\end{urlList}


\section{Psychosoziale Beratung}

Wenn du das Gefühl hast, die Kontrolle zu verlieren, oder nicht mehr mit
dem Studium und/oder den Menschen um dich herum zurecht kommst, wende dich an die kostenloase Psychosoziale Beratung des Studentenwerks.

\begin{urlList}
	\urlItem{http://www.studentenwerk-muenchen.de/beratungsnetzwerk/psychosoziale-und-psychotherapeutische-beratung/}
\end{urlList}


\section{Weitere Beratung des Studentenwerks}
Helene-Mayer-Ring 9 (U3 Olympiazentrum)

\begin{itemize}
	\item Allgemeine und Soziale Beratung
	\item Psychotherapeutische Beratungsstelle
	\item Studienkreditberatung
	\item Rechtsberatung
	\item Wohnungsberatung/\,Privatzimmervermittlung
	\item Beratungsstelle „Sexuelle Belästigung, Diskriminierung und Gewalt“
	\item Beratung für ausländische Studierende
\end{itemize}

\begin{urlList}
	\urlItem{http://studentenwerk-muenchen.de/beratungsnetzwerk}
\end{urlList}


\section{Nightline München}

Die Nightline München ist ein Zuhörtelefon von Studika für Studika,
das abends und nachts zu erreichen ist. Am Telefon sitzen ehrenamtlich
tätige Studika, die dir mit einem offenen Ohr beistehen.

\begin{urlList}
	\urlItem{http://www.nightline.mhn.de/}
\end{urlList}


\section{Kirchliche Beratung}
Die christlichen Hochschulgemeinden bieten neben ihrem konfessionellen Angebot auch überkonfessionelle und psychologische Beratung und Aktivitäten, wie Ausflüge, Workshops, Spieleabende\ldots

\begin{urlList}
	\urlItem{http://www.khg.lmu.de}
	\urlItem{http://www.esg.lmu.de}
\end{urlList}
